% GENERATED FILE. Edit canonical lessons, then run npm run generate:books.
% canonical-source-hash: fnv1a64:749cef456353a39e
% canonical-lessons: FR-C107-supermarche, FR-C107-cinema, FR-C107-stade, FR-C107-piscine, FR-C107-epicerie

\chapter{Supermarket, Cinema, Stadium, Swimming pool, Grocer's shop}
\label{ch:fr-supermarket-cinema-stadium-swimming-pool-grocer-s-shop}
\hlchaptermodality{107}

\begin{chapteropening}
\textbf{By the end of this chapter:} \emph{I can name a supermarket, a cinema, a stadium, a swimming pool, a grocer's shop.}

Complete the last lesson of chapter 107: I can name a supermarket, a cinema, a stadium, a swimming pool, a grocer's shop.
\end{chapteropening}

\section[le supermarché]{le supermarché — a supermarket}

\label{lesson:FR-C107-supermarche}



\begin{quote}
Before the new one: say the French for a car park, then the French for a lift.
\end{quote}

\subsection*{You'll want to know: le supermarché}
\textbf{le supermarché} — \textquotedblleft{}a supermarket\textquotedblright{}.

\begin{cousinweb}
\textbf{supermarché} — built from \emph{super-}, Latin for \textquotedblleft{}above\textquotedblright{}, and the French word for \textquotedblleft{}market\textquotedblright{}, which $\leftarrow$ Latin \emph{mercatus}, \textquotedblleft{}trade\textquotedblright{} (Lewis \& Short); modelled on English \emph{supermarket}.
\end{cousinweb}

\begin{grammarlens}[title={What you've built}]
One more word for this chapter.
\end{grammarlens}

\subsection*{Your turn}
\begin{itemize}
\raggedright
  \item \emph{Say it:} \emph{le supermarché}
  \item \emph{Say it:} \emph{le supermarché}, once more
  \item \emph{Say it:} say \emph{le supermarché}
  \item \emph{Recall:} say the French for a car park, then the French for a lift, then say \emph{le supermarché} again
\end{itemize}

\begin{tcolorbox}[breakable,colback=teal!4,colframe=teal!35!black,title={Before you move on}]
What does le supermarché mean? (\textquotedblleft{}A supermarket\textquotedblright{}.) Say it once more.
\end{tcolorbox}
\section[le cinéma]{le cinéma — a cinema}

\label{lesson:FR-C107-cinema}



\begin{quote}
Before the new one: say the French for a lift, then the French for a supermarket.
\end{quote}

\subsection*{You'll want to know: le cinéma}
\textbf{le cinéma} — \textquotedblleft{}a cinema\textquotedblright{}.

\begin{cousinweb}
\textbf{cinéma} — short for \emph{cinématographe}, a coinage of the 1890s from Greek \emph{kinema}, \textquotedblleft{}movement\textquotedblright{}, and \emph{graphein}, \textquotedblleft{}to write\textquotedblright{}. (EtymDB 2.1)
\end{cousinweb}

\begin{grammarlens}[title={What you've built}]
One more word for this chapter.
\end{grammarlens}

\subsection*{Your turn}
\begin{itemize}
\raggedright
  \item \emph{Say it:} \emph{le cinéma}
  \item \emph{Say it:} \emph{le cinéma}, once more
  \item \emph{Say it:} say \emph{le cinéma}
  \item \emph{Recall:} say the French for a lift, then the French for a supermarket, then say \emph{le cinéma} again
\end{itemize}

\begin{tcolorbox}[breakable,colback=teal!4,colframe=teal!35!black,title={Before you move on}]
What does le cinéma mean? (\textquotedblleft{}A cinema\textquotedblright{}.) Say it once more.
\end{tcolorbox}
\section[le stade]{le stade — a stadium}

\label{lesson:FR-C107-stade}



\begin{quote}
Before the new one: say the French for a supermarket, then the French for a cinema.
\end{quote}

\subsection*{You'll want to know: le stade}
\textbf{le stade} — \textquotedblleft{}a stadium\textquotedblright{}.

\begin{cousinweb}
\textbf{stade} $\leftarrow$ Latin \emph{stadium}, from Greek \emph{stadion}, a measure of length and the racecourse of that length — a learned borrowing. (EtymDB 2.1)
\end{cousinweb}

\begin{grammarlens}[title={What you've built}]
One more word for this chapter.
\end{grammarlens}

\subsection*{Your turn}
\begin{itemize}
\raggedright
  \item \emph{Say it:} \emph{le stade}
  \item \emph{Say it:} \emph{le stade}, once more
  \item \emph{Say it:} say \emph{le stade}
  \item \emph{Recall:} say the French for a supermarket, then the French for a cinema, then say \emph{le stade} again
\end{itemize}

\begin{tcolorbox}[breakable,colback=teal!4,colframe=teal!35!black,title={Before you move on}]
What does le stade mean? (\textquotedblleft{}A stadium\textquotedblright{}.) Say it once more.
\end{tcolorbox}
\section[la piscine]{la piscine — a swimming pool}

\label{lesson:FR-C107-piscine}



\begin{quote}
Before the new one: say the French for a cinema, then the French for a stadium.
\end{quote}

\subsection*{You'll want to know: la piscine}
\textbf{la piscine} — \textquotedblleft{}a swimming pool\textquotedblright{}.

\begin{cousinweb}
\textbf{piscine} $\leftarrow$ Latin \emph{piscina}, \textquotedblleft{}a fishpond\textquotedblright{}, from \emph{piscis}, \textquotedblleft{}fish\textquotedblright{} (EtymDB 2.1) — a learned borrowing: the same Latin fish as \emph{poisson}.
\end{cousinweb}

\begin{grammarlens}[title={What you've built}]
One more word for this chapter.
\end{grammarlens}

\subsection*{Your turn}
\begin{itemize}
\raggedright
  \item \emph{Say it:} \emph{la piscine}
  \item \emph{Say it:} \emph{la piscine}, once more
  \item \emph{Say it:} say \emph{la piscine}
  \item \emph{Recall:} say the French for a cinema, then the French for a stadium, then say \emph{la piscine} again
\end{itemize}

\begin{tcolorbox}[breakable,colback=teal!4,colframe=teal!35!black,title={Before you move on}]
What does la piscine mean? (\textquotedblleft{}A swimming pool\textquotedblright{}.) Say it once more.
\end{tcolorbox}
\section[l'épicerie]{l'épicerie — a grocer's shop}

\label{lesson:FR-C107-epicerie}



\begin{quote}
Before the new one: say the French for a stadium, then the French for a swimming pool.
\end{quote}

\subsection*{You'll want to know: l'épicerie}
\textbf{l'épicerie} — \textquotedblleft{}a grocer's shop\textquotedblright{}.

\begin{cousinweb}
\textbf{épicerie} — built from the French word for \textquotedblleft{}spice\textquotedblright{}, with \emph{-erie}, \textquotedblleft{}a shop for\textquotedblright{}. The spice word $\leftarrow$ Latin \emph{species}, \textquotedblleft{}a kind; wares\textquotedblright{} (Lewis \& Short).
\end{cousinweb}

\begin{grammarlens}[title={What you've built}]
One more word for this chapter.
\end{grammarlens}

\subsection*{Your turn}
\begin{itemize}
\raggedright
  \item \emph{Say it:} \emph{l'épicerie}
  \item \emph{Say it:} \emph{l'épicerie}, once more
  \item \emph{Say it:} say \emph{l'épicerie}
  \item \emph{Recall:} say the French for a stadium, then the French for a swimming pool, then say \emph{l'épicerie} again
\end{itemize}

\begin{tcolorbox}[breakable,colback=teal!4,colframe=teal!35!black,title={Before you move on}]
What does l'épicerie mean? (\textquotedblleft{}A grocer's shop\textquotedblright{}.) Say it once more.
\end{tcolorbox}
